\section{Evidência de IMP-NOTIF-002 (marcar como lida, M9 + server-owned)}

\textbf{Feature:} \texttt{IMP-NOTIF-002} --- marcar uma notificação como lida.

\textbf{Estado anterior:} \texttt{DIVERGENTE}. A callable usava apenas
\texttt{request.auth.uid}, sem reler atividade/versão/papel persistidos, e as
Rules permitiam \texttt{update}/\texttt{delete} diretos do dono, contrariando o
modelo server-owned e ``marcar lida só servidor''.

\textbf{Estado depois:} \texttt{NÃO DIVERGENTE} para o escopo examinado.

\subsection{Contrato}

Fontes: Seção~7 \texttt{RN-M13-02} e a matriz de autorização M9 (destinatário lê
a própria notificação; marcar lida só pelo servidor). A autoridade é a base M9
persistida (usuário ativo, versão corrente, papel persistido) e o ownership do
caminho (\texttt{Usuarios/\{uid\}/Notificacoes}).

\subsection{Mudança}

\begin{itemize}
  \item \path{functions/src/notificacoes.ts}: \texttt{marcarNotificacaoComoLida}
  passou a extrair claims e a reexecutar
  \texttt{resolverAutoridadePersistidaTx(tx, claims, PAPEIS\_CONHECIDOS)} dentro
  da transação; a notificação continua acessível apenas ao próprio UID;
  \item \path{firestore.rules}: a subcoleção \texttt{Notificacoes} passa a
  \texttt{allow read: if isOwner(uid); allow write: if false;}, tornando a
  marcação exclusivamente server-owned.
\end{itemize}

\subsection{TDD}

\small
\begin{longtable}{@{}p{0.30\textwidth}p{0.52\textwidth}@{}}
\toprule
ID & Resultado \\
\midrule
\endhead
TEST-INT-NOTIF-M9-001 & Dono ativo e corrente marca; \texttt{lida\_em} definido \\
TEST-INT-NOTIF-M9-002 & Não autenticado nega \\
TEST-INT-NOTIF-M9-003 & Usuário inativo nega e não altera \\
TEST-INT-NOTIF-M9-004 & Token obsoleto nega \\
TEST-INT-NOTIF-M9-005 & Claim sem papel persistido nega \\
TEST-INT-NOTIF-M9-006 & Repetição preserva \texttt{lida\_em} \\
TEST-INT-NOTIF-M9-007 & Notificação inexistente nega \\
TEST-RULES-NOTIF-001 & Dono com autoridade lê a própria notificação \\
TEST-RULES-NOTIF-002 & \texttt{update} direto negado \\
TEST-RULES-NOTIF-003 & \texttt{delete} direto negado \\
TEST-RULES-NOTIF-004 & \texttt{create} direto negado \\
TEST-RULES-NOTIF-005 & Terceiro não lê notificação alheia \\
TEST-RULES-NOTIF-006 & Token obsoleto não lê \\
TEST-RULES-NOTIF-007 & Usuário inativo não lê \\
\bottomrule
\end{longtable}
\normalsize

RED: três casos de integração (inativo, versão obsoleta, sem papel) e dois de
Rules (\texttt{update}/\texttt{delete}) falharam contra a implementação/Rules
anteriores. GREEN sob Node 20.19.6: integração 7/7, Rules 69/69.

\subsection{Regressão}

Sob Node 20.19.6: \texttt{build} PASS; \texttt{test:unit} PASS (71/71);
integração completa 117 passam / 14 falham (14 legadas em reagentes/patrimônio/
roteiros); Rules 69/69.

\subsection{Hardening de consistência persistida (M13)}

Auditoria independente apontou que o fechamento anterior provava a autoridade M9
e o server-owned, mas não a consistência da entidade formal \texttt{\#M13Notificacao}
no momento da marcação. Duas lacunas foram corrigidas.

\paragraph{Lacuna 1 --- UID do caminho vs.\ \texttt{id\_destinatario}.}
A norma (RN-M13-01; CUE \texttt{\#M13Notificacao} e \texttt{\#M13Marcacao};
Alloy \texttt{MarcarSoProprias}) exige que o \texttt{uid} do caminho e
\texttt{id\_destinatario} sejam o destinatário autenticado. Um documento em
\texttt{Usuarios/u1/...} com \texttt{id\_destinatario=u2} não pode ser marcado
por \texttt{u1}. \texttt{marcarNotificacaoComoLida} passou a exigir
\texttt{id\_destinatario == claims.uid} e falha fechado. As Rules de leitura
passaram a exigir \texttt{resource.data.id\_destinatario == uid} (CUE
\texttt{\#M13Leitura}: \texttt{id\_destinatario != uid} $\Rightarrow$
\texttt{autorizado=false}), fechando a leitura de documento endereçado a
terceiro ainda que situado no próprio caminho.

\paragraph{Lacuna 2 --- coerência \texttt{lida}/\texttt{lida\_em}.}
CUE \texttt{\#M13Notificacao} estabelece \texttt{lida=false} $\Rightarrow$
\texttt{lida\_em=null} e \texttt{lida=true} $\Rightarrow$ \texttt{lida\_em}
presente; a Alloy impõe a mesma coerência. Não há mecanismo normativo de
reparo. Estado persistido incoerente (\texttt{lida=true, lida\_em=null} ou o
inverso) faz a marcação falhar fechado com \texttt{failed-precondition}, sem
mutação e sem auto-reparo silencioso.

\small
\begin{longtable}{@{}p{0.30\textwidth}p{0.52\textwidth}@{}}
\toprule
ID & Resultado \\
\midrule
\endhead
TEST-INT-NOTIF-M13-008 & Caminho de \texttt{u1} com \texttt{id\_destinatario=u2} nega e não altera \\
TEST-INT-NOTIF-M13-009 & \texttt{lida=true, lida\_em=null} falha fechado sem mutação \\
TEST-INT-NOTIF-M13-010 & \texttt{lida=false, lida\_em} preenchido falha fechado sem mutação \\
TEST-INT-NOTIF-M13-011 & Notificação sem \texttt{id\_destinatario} nega (fail-closed) \\
TEST-RULES-NOTIF-008 & Dono não lê documento com \texttt{id\_destinatario} divergente \\
TEST-RULES-NOTIF-009 & Leitura coerente (\texttt{id\_destinatario == uid}) é permitida \\
\bottomrule
\end{longtable}
\normalsize

RED: os casos 008--011 falharam contra a implementação anterior (marcação
aceitava documento alheio e estado incoerente) e o caso de Rules 008 falhou
contra as Rules anteriores. GREEN sob Node 20.19.6: integração de notificações
11/11 e Rules de notificações 9/9. Regressão: \texttt{build} PASS; unit 71/71;
Rules 75/75; integração completa 129 passam / 14 falham (falhas legadas em
reagentes/patrimônio/roteiros, sem relação com a correção). Documentação
matricial: achados de ``Achados de maior impacto'' atualizados para o estado
corrente (M7 canônico existe e cobre papéis; notificações fechadas).

\textbf{Estado final:} \texttt{IMP-NOTIF-002} = \texttt{NÃO DIVERGENTE} no
escopo auditado. As demais linhas de notificação permanecem abertas.
